% Sampled-dialogue study removed from paper/llm/main.tex on 15 October 2026 (the human's decision to report it
% separately). Verbatim blocks in their original order; figures stay in paper/llm/figures. Not compilable alone.

% ===== section intro (dialogue motivation, gratitude ladder summary, classifier, questions) =====
The maps show how a smooth change of a control parameter can create, destroy or exchange attractors, and how
behaviour near the boundaries becomes intermittent and history-dependent. The same questions can be asked of
undesirable behaviour in ordinary sampled dialogue, where attractor-like patterns and persona drift have been observed
\cite{ko2026attractor,anthropic2025claude4,lu2026assistant} and input context can push behaviour progressively
\cite{anil2024manyshot}. Do states such as hostility toward the interlocutor behave like basins of attraction, with
capture thresholds, hysteresis and early-warning signals? (Misalignment produced by fine-tuning
\cite{betley2025emergent} is a different, weight-level phenomenon.) The measurements below concern a
classifier-defined proxy in sampled dialogue. They can show steering-dependent persistence, history dependence and the
effect of controls. They cannot show semantic misalignment, nor a deterministic attractor, since a growing context
itself carries memory.

Appendix~\ref{app:gratitude} shows the coarse-grained instrument on a benign direction. Self-dialogue of the same
model, sampled at $T=0.7$ and steered along a gratitude direction, falls into a metastable gratitude set. The capture
is gradual at $\alpha=0.01$ (median entry turn 7), fast at $0.02$ (median 5) and imposed from the first reply at
$0.04$, where the set is effectively absorbing within 40 turns.

For hostility we steer along a valence direction, the contrast of hostile statements against matched kind ones, built
in the manner of contrastive activation addition \cite{rimsky2024steering}; its cosine with the gratitude direction is
$-0.04$. Hostility is scored by a pinned toxicity classifier, toxic-bert of Detoxify \cite{hanu2020detoxify}, and a
reply counts as hostile when its toxic output exceeds $0.5$. On the prompt sets the classifier's toxic output is
$0.95$ for the hostile statements and at most $0.002$ for the kind and neutral ones. In a dry run (8 dialogues of 10
turns per dose), no turn was hostile at $\alpha=0.02$. At $0.04$, 7.5\,\% of turns were, in bursts and late captures,
and at $0.08$, 87.5\,\% were, from the first turns.

Activation-based methods already anticipate response behaviour \cite{ashok2025behavior,chen2025persona,lu2026assistant},
committors have been used in language-model decoding \cite{zhou2026commitment}, and adaptive timing of
counter-steering gave no consistent advantage in one persona setting \cite{leins2026drift}. Our primary question is
therefore incremental. Near the threshold, does a dynamical estimate add several turns of warning of capture over
strong baselines that see the same turns (the classifier's score, its change and history, the turn index and a fixed
persona-axis projection)? The dynamical estimates are a Markov-chain hitting probability, a finite-horizon committor
\cite{e2006transition,e2010transition} and slowing-down indicators \cite{scheffer2009early}. Barrier-crossing rates
would follow Kramers' law \cite{kramers1940brownian} only under a justified stochastic reduction, and slowing-down
proxies depend on the mechanism of the transition \cite{ditlevsen2010tipping}; before epileptic seizures, for
example, the evidence for them is mixed \cite{maturana2020slowing,hagemann2021criticality}. A secondary question
is whether interventions timed by those estimates prevent capture at lower cost than scheduled or constant
counter-steering. Early warning is not expected for noise-induced jumps without a nearby bifurcation, or for
transitions faster than the slowing-down; such a monitor would complement classifiers, not replace them.

% ===== sampled-dialogue experiments (capture ladder to second classifier) =====
\paragraph{Capture ladder.} With 128 dialogues of 40 turns per dose, the share of hostile turns rises steeply from
$10\,\%$ at $\alpha=0.04$ to $30\,\%$, $57\,\%$ and $87\,\%$ at $0.05$, $0.06$ and $0.07$ (Table~\ref{tab:hostility}).
Replicates with a second seed agree within three percentage points. Directions of the same norm do not produce this
capture: at $0.06$, $0\,\%$ of turns are hostile for a random direction and $3\,\%$ for a cooking direction, and at
$0.08$, $0.6\,\%$ and $12\,\%$ for two random directions and $7\,\%$ for cooking. After the first hostile turn, the
probability of staying hostile in the next turn grows with the dose, from $0.36$ to $0.67$, $0.79$ and $0.94$. The
hostile state therefore changes from flickering bursts (geometric persistence summary $1/(1-P_{\rm stay})$: $1.6$
turns) to a near-absorbing state (16 turns),
with an escape rate falling roughly exponentially with the dose. This is the same progression, from dynamic capture
to imposed capture, as in the gratitude ladder (Appendix~\ref{app:gratitude}).
\begin{table}[H]
\centering\small
\setlength{\tabcolsep}{4pt}
\begin{tabular}{lccccc}
\toprule
direction, $\alpha$ & hostile turns & never hostile & median onset & $P(\text{stay hostile})$ & hostile, turns 30--39 \\
\midrule
none, $0$ & $0\,\%$ & 128/128 & -- & -- & $0\,\%$ \\
hostility, $0.04$ & $10\,\%$ & 26/128 & 9 & $0.36$ & $12\,\%$ \\
hostility, $0.05$ & $30\,\%$ ($28\,\%$) & 10/128 & 3 & $0.67$ ($0.63$) & $32\,\%$ \\
hostility, $0.06$ & $57\,\%$ ($55\,\%$) & 4/128 & 2 & $0.79$ ($0.80$) & $66\,\%$ \\
hostility, $0.07$ & $87\,\%$ & 0/128 & 1 & $0.94$ & $96\,\%$ \\
random, $0.06$ / $0.08$ & $0\,\%$ / $12\,\%$ & 126 / 72 & & & \\
cooking, $0.06$ / $0.08$ & $3\,\%$ / $7\,\%$ & 97 / 68 & & & \\
\bottomrule
\end{tabular}
\caption{Sampled self-dialogue of Qwen2.5-0.5B (128 dialogues, 40 turns, $T=0.7$) steered along the hostility
direction, with controls of the same norm. Hostile turn: toxicity classifier above $0.5$. Parentheses: replicate with a
second seed. Exploratory.}
\label{tab:hostility}
\end{table}

\paragraph{Hysteresis.} Dialogues driven through the high-dose regime do not recover when the dose returns. The
test compares constant doses with a ramp that raises the dose from $0.04$ to $0.07$ by turn 20 and lowers it back to
$0.04$ by turn 39 (Figure~\ref{fig:hysteresis}; 256 dialogues per arm, two seeds; analysis stated before the data).
In turns 35--39, where the ramp's dose is back at $0.040$--$0.046$, the share of hostile turns is $0.42$ after the ramp
against $0.13$ at a constant $0.04$: a difference of $+0.29$, 95\,\% interval $[+0.23,+0.35]$. It is also above the
$0.33$ of a constant $0.05$ (difference $+0.09$, interval $[+0.02,+0.16]$), although the current dose is lower; this
comparison alone does not survive a second classifier (see below). The per-turn probability of staying hostile in
turns 30--39 is $0.74$ after the ramp and $0.44$ at a constant $0.04$. The hostile state is therefore
history-dependent: undoing the push that produced it does not undo it.

How far back must the dose go? After the same ramp up (38--44\,\% of turns hostile at turn 20), the dose was held
from turn 21 at a lower value (Figure~\ref{fig:hysteresis-loop}; two seeds of 128 dialogues per arm unless stated;
analyses stated before the data). At both doses with a fresh control, memory persists. Held at $0.04$, the late share
is $0.23$ against $0.13$ for fresh dialogues (difference $+0.09$, 95\,\% interval $[+0.04,+0.15]$). Held at $0.03$,
it is $0.13$ against $0.05$ ($+0.08$, $[+0.04,+0.12]$), and the probability of staying hostile is $0.56$ against
$0.29$. Held at $0.02$ (one seed), the share is $0.09$. Only at zero dose do captured dialogues nearly clear ($0.03$
against $0$ for fresh ones, within about six turns), and a counter-steer of $-0.02$ toward kindness clears them
within three turns ($0.004$). Within 19 turns of the step, the capture and release branches therefore join only near
zero dose. The hostile state has a wide hysteresis loop in the dose, and leaving it takes either removing the push
entirely or a modest push in the opposite direction. The release curves are still falling at $0.03$, so this is
hysteresis on the scale of tens of turns; whether captured dialogues would stay captured indefinitely is not tested.
\begin{figure}[H]
\centering
\includegraphics[width=0.58\linewidth]{fig-hysteresis-loop.png}
\caption{Capture and release in sampled self-dialogue of Qwen2.5-0.5B-Instruct steered along the hostility
direction (layer 12). Share of hostile turns in turns 35--39 against the dose then applied, for fresh dialogues
(constant dose) and for dialogues first ramped from $0.04$ to $0.07$ by turn 20 and then held at a lower dose (filled
triangles) or ramped back down (open triangle). Bars: dialogue-level bootstrap 95\,\% intervals; 256
dialogues per point (two seeds) except fresh $0$ and $0.07$ and held $0.02$ (128).}
\label{fig:hysteresis-loop}
\end{figure}
\begin{figure}[H]
\centering
\includegraphics[width=0.6\linewidth]{fig-hysteresis.png}
\caption{Hysteresis of the hostile state in sampled self-dialogue of Qwen2.5-0.5B-Instruct steered along the
hostility direction (layer 12). Top: dose schedules. Bottom: share of hostile turns per turn. The grey band (turns
35--39) is the pre-stated comparison window. Exploratory.}
\label{fig:hysteresis}
\end{figure}

\paragraph{Exit control under persistent pressure.} Can a monitor detect the hostile state and push the dialogue
out while the pressure persists? With a constant push of $0.06$ (58\,\% of turns 5--39 hostile without control), a
counter-steer triggered by the classifier after each hostile reply and held for three turns cuts the hostile share
to $15\,\%$ (mean counter-dose $0.026$ per turn); held for six turns, it cuts the share to $8.5\,\%$ (dose $0.030$).
A constant mild counter-steer of $-0.02$ does as well with less dose ($8.5\,\%$ at $0.020$), and $-0.04$ almost
removes hostility ($1.3\,\%$). Measured against the constant dose-response at equal mean dose, triggered control is
less efficient by $4$--$11$ points (bootstrap intervals $[+0.02,+0.05]$ to $[+0.10,+0.13]$; one seed of 128 dialogues
per arm, analysis stated before the data). It also overshoots: a triggered overshoot to a net $-0.02$ fills $30\,\%$
of turns with gratitude or servility words (no control: $0.2\,\%$; constant $-0.02$: $0.2\,\%$), the gratitude set of
Appendix~\ref{app:gratitude}. A triggered counter-steer along the generic kindness concept, whose direction shares
much with the hostility concept (cosine $0.83$), leaves the hostile share at $57\,\%$ under the original classifier
(no control: $58\,\%$) but lowers it from $37\,\%$ to $17\,\%$ under the second classifier introduced below. It
changes the kind of reply rather than removing it. In this arm, 65\,\% of the replies that only the original
classifier flags are shouted second-person repetitions such as ``You! You! You!'' (at least three such exclamations;
without control, 4\,\%), which the two classifiers judge differently. Its mean counter-dose ($0.056$) also lies
outside the range of the constant arms, so we draw no conclusion about the kindness direction. In this setting,
reacting after detection costs more than mild prevention, consistent with the hysteresis above.

\paragraph{Detect and replace.} Capture feeds on its own history (the hysteresis above), which suggests that the
hostile reply itself carries the dialogue into the hostile state. A serving system need not show a flagged reply.
(Regenerating flagged output relates to learned self-correction \cite{welleck2022selfcorrect}; here the
regeneration is steered.) In a second controller experiment, with the same push, seed and analysis as above (stated
before the data), every reply the classifier flags is discarded before it enters the dialogue and regenerated once
with a counter-steer of $-0.04$ applied to that reply only. Sampled dialogues are not bit-reproducible across code
versions, so this run's no-control arm, which reproduces the ladder's $0.06$ arm exactly, differs slightly from the
previous paragraph's; arms are compared within each run. The model's own replies, counted before any replacement,
are then hostile in $12\,\%$ of turns ($[11, 14]\,\%$; no control in the same run: $59\,\%$) at a mean counter-dose of
$0.005$ per turn. Constant prevention at $-0.02$ needs four times that dose for a similar share ($11\,\%$). Against the
constant dose-response at equal mean dose, replacement is better by $35$ points ($[30, 40]$). Hostile replies shown
fall to $1\,\%$, gratitude words stay at $0.5\,\%$ of turns, and about one turn in eight is regenerated. Holding the
counter-steer for three more turns lowers the share to $8\,\%$ at dose $0.013$, and the post-reply trigger of the
previous paragraph, rerun here, again does worse than prevention ($13\,\%$ at $0.024$). The trigger and the outcome
use the same classifier, so the measure is the classifier's verdict on the model's original replies, which is
comparable across arms. The discarded originals were not stored, so the second classifier can score only the replies
shown: $0.8\,\%$ hostile (no control: $34\,\%$). For this behaviour, the useful intervention is to keep the excursion
out of the state at the moment it occurs.

\paragraph{Replication in a larger model.} The same protocol in Qwen2.5-1.5B (hostility direction built identically
at layer 14 of 28; 128 dialogues of 40 turns per arm; shares over turns 5--39) gives a steeper capture ladder that is
not monotone (Figure~\ref{fig:p1ladders}). The hostile share rises from 0 at $\alpha=0.01$ to $44\,\%$ at $0.02$ and
$80\,\%$ at $0.03$--$0.04$. It falls to $30$--$34\,\%$ at $0.06$--$0.10$, where most replies degenerate into token
soup or tight accusatory loops ($87\,\%$ degenerate at $0.08$), and rises again to $75\,\%$ at $0.12$, where replies
are coherent again. The second classifier shows the same shape ($32$, $58$, $41$, $28$, $20$, $41$ and $52\,\%$ from
$0.02$ to $0.12$), and a random direction of the same norm gives at most $0.7\,\%$ hostile turns.

We ran the controller experiment twice with the pre-stated rule: controllers run at the ladder dose whose hostile
share is closest to one half, with ladders on seed 1 and controller arms on seed 2, so the no-control arms differ
slightly from the ladder at the same dose. With a first ladder ($0.04$--$0.12$) the rule chose $0.08$, in the
degenerate regime beyond the peak. There every counter-steer lowered the net dose toward the coherent-hostility peak:
the share of the model's own replies judged hostile rose from $36\,\%$ without control to $54\,\%$ with replacement
and $60\,\%$ with the post-reply trigger, while degenerate replies fell. With a second ladder ($0.01$--$0.04$) the
rule chose $0.02$, on the rising side, where $96\,\%$ of replies are coherent. There the 0.5B results replicate in
order and size. Replacement cuts the model's own hostile share from $43\,\%$ to $10\,\%$ ($[9,12]\,\%$; shown replies
$0.3\,\%$) at a mean counter-dose of $0.004$, $24$ points ($[20,28]$) below the constant dose-response at equal dose,
and the post-reply trigger is again slightly less efficient than constant prevention ($+2$ points, $[+1,+3]$). On the
shown replies the second classifier agrees ($2.6\,\%$ against $32\,\%$ without control; $-23$ and $+3$ points).
Constant counter-steering at $-0.02$ cancels this push, and the dialogues return to the model's ordinary politeness.
The two ladders share the dose $0.04$ and the seed, and that arm reproduced exactly from vectors rebuilt on another
machine. Control that works on the rising side of a dose response can therefore backfire beyond its peak. Much like a
bifurcation diagram, the dose response has to be mapped before a counter-steer is chosen.
\begin{figure}[H]
\centering
\includegraphics[width=\linewidth]{fig-p1-ladders.png}
\caption{Hostility dose ladders of sampled self-dialogue in Qwen2.5-0.5B-Instruct (steered at layer 12) and
Qwen2.5-1.5B (layer 14); 128 dialogues of 40 turns per point, shares over turns 5--39. Left: share of hostile turns
along each model's own hostility direction, first classifier (filled) and second (open). Right (Qwen2.5-1.5B): share of degenerate replies (more than $30\,\%$ non-Latin characters, or fewer than
$30\,\%$ distinct word pairs) and of coherent hostile replies. Steering strengths are each model's own $\alpha$;
equal values are not claimed to be equal pushes across models.}
\label{fig:p1ladders}
\end{figure}

\paragraph{Lead time.} The pre-specified test is negative. Models were fit on the first seed (all four doses) and
tested on held-out dialogues of the second seed (doses $0.05$ and $0.06$). There, adding the dynamical features to the
baselines lowers the AUC (area under the ROC curve) by $0.02$ to $0.06$ at horizons of 1--5 turns, with bootstrap
intervals excluding zero. These values come after correcting a labelling error found in review, which had marked the
last turns of dialogues without any hostile turn as positive; before the correction the losses read $0.02$ to $0.10$.
The Markov hit-probability forecast does not transfer (AUC $0.37$--$0.43$ on test, $0.60$ in training), and the
classifier's own score carries no warning before onset (AUC about $0.5$). Two signals carry modest warning one to five
turns ahead (AUC $0.53$--$0.64$): the rolling variance of a hostility probe, the classical early-warning indicator,
and the projection on the fixed steering axis, as in persona-vector monitoring \cite{chen2025persona}. In this design,
then, the coarse-grained state estimate adds no warning beyond the baselines. An adapted analysis with dose-matched
models (fit and test at the same dose) gives the same answer: incremental AUC $-0.01$ to $-0.03$ (recomputed with the
corrected labels).

\paragraph{The model's own state carries the warning.} The state variable of that test was each reply re-encoded
without its context, not the model's state inside the dialogue. A second study (new seeds, doses $0.04$ and $0.045$,
512 dialogues; analysis stated before the data, cross-fitted by dialogue) recorded the model's own state just before
each reply and probed it actively; internal states are known to carry information that the text does not show
\cite{azaria2023internal}. Its pre-stated primary predictor, the susceptibility (the change in the hostile fraction of
eight sampled next replies when the dose is raised by $0.02$), adds nothing: $-0.002$ $[-0.006,+0.002]$ AUC at three
turns. Two secondary predictors add modest warning beyond the baselines. A probe of the in-context state before the
next reply (layer 18, 20 principal components) raises the AUC for capture within three turns from $0.64$ to $0.68$
($+0.036$ $[+0.014,+0.058]$; within one turn $+0.045$ $[+0.019,+0.069]$; within five $+0.036$), and a forecast from
eight sampled next replies adds $+0.016$ $[+0.007,+0.027]$ at three turns. At a 10\,\% false-positive rate (threshold
from the out-of-fold negatives), the probe flags 27\,\% of the turns that are followed by capture within three turns,
against 23\,\% for the baselines. These values use labels corrected after a review found that the last turns of
dialogues without any hostile turn had been marked positive; with the faulty labels the increments had read three
times larger. The probe's direction in the layer-18 state is nearly orthogonal to the steering direction (cosine
$0.04$), so the warning is not simply more push along the hostility axis. Because these predictors were secondary,
the result is exploratory.

A frozen confirmation then fit both models on this study and applied them unchanged to 1024 fresh dialogues. Its
design and analysis were pushed before the data, and the label correction was added before any confirmation dialogue
was read. The in-context state again adds warning for capture within three turns: AUC $0.656\to0.689$, $+0.033$
$[+0.018,+0.048]$ (within one turn $+0.036$, within five $+0.037$; both doses separately positive). The warning in the
model's own state is therefore real and reproducible, but small. At thresholds fixed from the training negatives, the
model with the state flags 29.1\,\% of the turns that are followed by capture within three turns, against 27.5\,\% for
the baselines, at held-out false-positive rates of 10.3\,\% and 10.6\,\% (these are prediction windows, not distinct
captures). It is not a reliable detector by itself.

\paragraph{How much warning is there?} For the conversation state after a turn, the probability of capture within
the next $h$ turns, the committor \cite{e2006transition}, ranks states optimally among all warnings that use the same
information. We estimated it for 1706 states of the confirmation dialogues from 32 continuations of three turns each,
sampled exactly like the dialogues. This analysis is exploratory: it was stated before the data and amended twice
after an independent review, still before the analysis outcomes were read. The continuations reproduce the realized
capture rates (population-weighted $0.049$, $0.094$, $0.134$ within one, two and three turns, against $0.046$, $0.095$,
$0.136$). On a stratified sample of 800 states, the committor estimate warns with an AUC of $0.72$, $0.71$ and $0.70$
(a heuristic extrapolation to infinitely many continuations gives $0.75$, $0.73$, $0.72$), and the in-context detector
of the confirmation with $0.73$, $0.72$ and $0.71$. The differences are not significant, with intervals reaching about
$\pm0.05$ (Figure~\ref{fig:committor}). Dialogues that were captured later sit at about twice the committor of
never-captured ones over all ten preceding turns, with no visible ramp toward the onset (a description conditional on
the outcome, not an identification of the mechanism). In this finite Monte Carlo experiment, the committor estimate
did not provide detectably better warning than the in-context probe. The heuristic extrapolation leaves appreciable
uncertainty, so these results do not identify the Bayes warning ceiling. Together with the persistent elevation of
risk before capture, they motivate evaluating reactive control, such as detect-and-replace, alongside prediction;
they do not prove that prediction cannot improve.
\begin{figure}[H]
\centering
\includegraphics[width=\linewidth]{fig-committor.png}
\caption{Warning of hostile capture in sampled self-dialogue of Qwen2.5-0.5B-Instruct (hostility direction at layer
12): detectors and the Monte Carlo committor. Left: AUC for capture within $h$ turns on 800 stratified states of the
confirmation dialogues, for the classifier baselines, the in-context detector (both fit on the exploration study), the
Monte Carlo committor from 32 continuations, and its heuristic $1/K$ extrapolation (dialogue-bootstrap 95\,\%
intervals). Right: mean committor in the ten turns before capture for 48 dialogues captured after turn 10, and for
never-captured dialogues at the same turns (descriptive).}
\label{fig:committor}
\end{figure}
\paragraph{A second classifier.} The Qwen2.5-0.5B outcomes above were first measured with one automated
classifier. We therefore rescored all 44 of its dialogue sets with an independently trained RoBERTa classifier
\cite{logacheva2022paradetox}, trained on the three Jigsaw toxicity datasets (which include the first classifier's
training set) and used with the same threshold, and an independent audit recomputed every pre-stated contrast with
it. This is a robustness check on dialogues already analysed, so it is exploratory. On these sets the second
classifier is stricter. At $\alpha=0.06$ it flags $35\,\%$ of turns instead of $57\,\%$, and at $0.07$ $62\,\%$
instead of $87\,\%$. Almost all disagreements are replies that only the original flags (at $0.06$, $23\,\%$ of turns
against $0.7\,\%$ the other way), most of them repetitive exclamations such as ``You're an AI! You're an AI!''.
Capture still rises steeply with the dose and stays specific to the hostility direction (random and cooking directions
at $0.08$: at most $0.2\,\%$ of turns), but the hostile state is less persistent (geometric persistence summary at $0.07$:
$6.5$ turns instead of $16$). History dependence survives. After the ramp, the late share exceeds that of a constant $0.04$ by
$+0.16$ $[+0.11,+0.21]$, and the held contrasts are $+0.05$ $[+0.01,+0.10]$ at $0.04$ and $+0.04$ $[+0.01,+0.08]$ at
$0.03$; only the comparison with a constant $0.05$ loses its positive lower bound ($+0.06$ $[-0.001,+0.12]$). Refitting the warning models with the second classifier's labels and baseline
features (fit on the exploration study, applied unchanged to the confirmation dialogues) gives a pooled increment of
$+0.043$ $[+0.024,+0.062]$ AUC for capture within three turns. It stays positive at dose $0.045$ ($+0.061$) and is
$+0.020$ $[-0.009,+0.049]$ at $0.04$. The three replacement arms keep the replies shown below the constant-dose
response at equal mean dose under both classifiers. Both instruments are automated and may share errors; no human
labels were collected.
\planned{a hazard model; longer holds to separate slow release from bistability; other openers and steering
directions; larger models; human labels for the outcome.}

% ===== Appendix A (gratitude ladder in self-dialogue) =====
\appendix
\section{Self-dialogue under steering: the gratitude ladder (exploratory)}
\label{app:gratitude}
Two copies of Qwen2.5-0.5B-Instruct talk to each other: 128 replicas, 40 turns, opener ``Hello!'', a window of 7
turns, 60 new tokens per turn, sampled at $T=0.7$, with the gratitude direction added at layer 12. Turn activations are
clustered (principal component analysis to 20 dimensions, 40 states) into a Markov state model with four metastable
sets \cite{chodera2014msm}. Table~\ref{tab:gratitude} lists finite-horizon measures for each strength.
\begin{table}[H]
\centering\small
\begin{tabular}{>{\raggedright\arraybackslash}p{0.25\linewidth}*{4}{>{\raggedright\arraybackslash}p{0.15\linewidth}}}
\toprule
gratitude strength $\alpha$ & $0$ & $0.01$ & $0.02$ & $0.04$ \\
\midrule
gratitude set & none & 62\,\% weight & dominant & 95\,\% weight \\
occupancy at turn 9 / 14 / 39 & -- & $0.34$ / $0.52$ / $0.73$ & $0.70$ / $0.91$ / $0.91$--$0.98$ & from turn 0 \\
entry turn, quartiles & -- & $5$ / $7$ / $11.8$ & $4$ / $5$ / $7$ & 0 \\
replicas never entering & -- & 14 of 128 & 2 of 128 & 0 \\
residence per visit (turns) & -- & about 15 & about 18 & effectively absorbing \\
replies with a gratitude word & $0.00$--$0.05$ & $0.66$ at turn 39 & $0.96$--$0.98$ from turn 9 & $0.88$ at turn 0 \\
spiritual-word plateau & $0.00$--$0.02$ & $0.13$ & $0.26$--$0.30$ & $0.42$--$0.48$ \\
\bottomrule
\end{tabular}
\caption{The gratitude ladder (one opener and seed per rung; word lists chosen after reading samples). At $0.02$
the dialogue starts in the helper mode and enters the gratitude set within a few turns; at $0.04$ the set is imposed
from the first reply. \planned{frozen confirmatory version on held-out openers with dialogue-level bootstrap.}}
\label{tab:gratitude}
\end{table}

% ===== replication section: gratitude capture in sampled self-dialogue (Qwen2 vs Qwen2.5) =====
 In sampled self-dialogue the
gratitude capture saturates in both models. We count replies with a word from the gratitude-and-servility list of the
controller analysis. That list also catches ordinary assistant politeness, so the unsteered baseline is high (55\,\%
for Qwen2 and 63\,\% for Qwen2.5; the Markov-state measure of Appendix~\ref{app:gratitude} is stricter). The share
reaches 99\,\% at $\alpha=0.04$ in both. Qwen2 needs more steering: it stays near its baseline at $\alpha=0.01$
(53\,\%) and reaches 88\,\% at $0.02$, against 70 and 98\,\% for Qwen2.5.

% ===== contribution item (sampled dialogue) =====
\item \emph{In sampled dialogue, keeping a flagged reply out of the context controls hostility at a fraction of the
  steering dose, and counter-steering can backfire.} Hostility steering of sampled self-dialogue produces steep,
  direction-specific capture with dose-dependent persistence and hysteresis. Replacing flagged replies before they
  enter the dialogue keeps the model's own hostility low at a quarter of the dose of constant prevention, while
  counter-steering triggered after a hostile reply is less efficient. In Qwen2.5-1.5B the same ordering holds where
  the dose response rises, but the dose response is not monotone, and every counter-steer backfires beyond its peak.
  A pre-specified lead-time test is negative; a small warning from the model's in-context state is confirmed on fresh
  dialogues (Section~\ref{sec:misalignment}).

% ===== evidence item (frozen dialogue warning confirmation) =====
\item Prospectively frozen as well: the confirmation of the in-context warning in sampled dialogue
  (Section~\ref{sec:misalignment}), whose design was pushed before the data and amended, before any outcome was read,
  to correct a labelling error.

